% =====================================================================
%  LAMPIRAN A — PERSYARATAN FISIK DAN TATA LETAK
%  Isi diambil langsung dari Panduan Skripsi Filkom UB v3.0 (wajib
%  dilampirkan verbatim di setiap skripsi/proposal, bukan hasil tulisan
%  penulis). Jangan diubah kecuali panduan resminya sendiri berubah.
% =====================================================================

\chapter*{Lampiran A Persyaratan Fisik dan Tata Letak}
\addcontentsline{toc}{chapter}{Lampiran A Persyaratan Fisik dan Tata Letak}

\subsection*{A.1 Kertas}

Kertas yang digunakan adalah HVS 70 gram berukuran A4. Apabila terdapat
gambar yang menggunakan kertas berukuran lebih besar dari A4, hendaknya
dilipat sesuai dengan aturan yang berlaku. Pengetikan hanya dilakukan
pada satu muka kertas, tidak bolak-balik.

\subsection*{A.2 Margin}

Batas pengetikan naskah adalah sebagai berikut:

\begin{itemize}
  \item Margin kiri: 4 cm
  \item Margin atas: 3 cm
  \item Margin kanan: 3 cm
  \item Margin bawah: 3 cm
\end{itemize}

\subsection*{A.3 Jenis dan Ukuran Huruf}

Jenis huruf yang dipakai adalah Calibri, dengan ketentuan sebagai
berikut:

\begin{itemize}
  \item Judul bab (level 1): 16 pt
  \item Judul subbab level 2: 14 pt
  \item Judul subbab level 3: 14 pt
  \item Judul subbab level 4: 12 pt
  \item Badan teks: 12 pt
\end{itemize}

Penggunaan jenis dan ukuran huruf ini harus konsisten di seluruh naskah.

\subsection*{A.4 Spasi}

Jarak standar antar baris dalam badan teks adalah satu spasi.

\subsection*{A.5 Kepala Bab dan Subbab}

Kepala bab terdiri dari kata ``BAB'' diikuti nomor bab dan judul bab,
misalnya ``BAB 1 PENDAHULUAN''. Kepala subbab diawali nomor sesuai
tingkat hierarkinya, diikuti judul subbab, misalnya ``1.2 Rumusan
Masalah''. Penomoran subbab disarankan tidak lebih dari 4 level
(maksimal X.X.X.X), tetapi sebaiknya hanya sampai 3 level. Kepala bab
dan subbab tidak boleh mengandung \emph{widow} atau \emph{orphan}
sehingga tampak menggantung atau terputus di awal/akhir halaman.

\subsection*{A.6 Nomor Halaman}

Bagian awal naskah menggunakan nomor halaman angka Romawi kecil (i, ii,
iii, dan seterusnya) dimulai dari sampul dalam. Bagian utama dan bagian
akhir menggunakan angka Arab (1, 2, 3, dan seterusnya) dimulai dari
Bab 1. Semua nomor halaman diletakkan di tengah bawah.
